% GENERATED FILE. Edit canonical lessons, then run npm run generate:books.
% canonical-source-hash: fnv1a64:89f0f96e0eefe482
% canonical-lessons: ES-C58-tampoco

\chapter{Tampoco --- The Same Frame, the Other Half}
\label{ch:tampoco}
\hlchaptermodality{220}

\begin{chapteropening}
\textbf{By the end of this chapter:} \emph{I can agree with a no as easily as with a yes.}

I can agree with a no as easily as with a yes.
\end{chapteropening}

\section[tampoco]{tampoco — the same frame, the other half}

\label{lesson:ES-C58-tampoco}



\begin{quote}
\emph{Yo también} agrees with somebody who said a \textbf{yes}.

But what do you say to somebody who said a \textbf{no}? English changes word there too — \emph{me too} becomes \emph{me neither}.
\end{quote}

\begin{sounds}
\begin{itemize}
\raggedright
  \item \texttt{k-hard} — the \emph{c} before \emph{o} is a plain \emph{k}. \textbf{tampoco} = \emph{tam-PO-ko}, with no accent to write, because the stress already falls where Spanish expects it.
\end{itemize}
\end{sounds}

\begin{cousinweb}
\textbf{tampoco} = \textbf{tan} + \textbf{poco}.

The same \emph{tan} you met yesterday. The new half is \emph{\textbf{poco}}, which means \emph{little} — a small amount of something.

\noindent
\begin{tabularx}{\linewidth}{@{}>{\raggedright\arraybackslash}X>{\raggedright\arraybackslash}X>{\raggedright\arraybackslash}X@{}}
\toprule
\textbf{} & \textbf{} & \textbf{} \\
\midrule
\emph{también} & tan + \textbf{bien} & so \textbf{well} \\
\emph{tampoco} & tan + \textbf{poco} & so \textbf{little} \\
\bottomrule
\end{tabularx}

One frame, two fillings. You are not learning a second word here so much as turning the first one over.
\end{cousinweb}

\begin{culture}
\begin{quote}
— \emph{No estoy cansado.} — \emph{\textbf{Yo tampoco.}}
\end{quote}

The rule is simple and it is about the \textbf{other} person's sentence, not yours. If what they said was positive, you answer \emph{yo también}. If it had a \emph{no} in it, you answer \emph{yo tampoco}.

Getting this wrong is one of the most audible mistakes a learner can make, and it is entirely avoidable: listen for the \emph{no}. If you heard one, use the \emph{poco} word.
\end{culture}

\subsection*{Your turn}
Say these aloud:

\begin{itemize}
\raggedright
  \item \emph{Yo también}, then \emph{Yo tampoco}
  \item \emph{No estoy cansado. — Yo tampoco.}
\end{itemize}

\emph{Twice through:}

\begin{tcolorbox}[breakable,colback=teal!4,colframe=teal!35!black,title={Before you move on}]
\emph{Tan} + \emph{bien} is \emph{so well}. What is \emph{tan} + \emph{poco}? (\textbf{So little} — \emph{\textbf{tampoco}}.)

And what in the other person's sentence decides between the two? (A \emph{\textbf{no}}.)
\end{tcolorbox}
